\subsection{基本群}

在下述命题中，\(f(\mathrm{G},\mathrm{T})\) 表示从 \(\varGamma(\mathrm{T})\) 到 \(\pi_{1}(\mathrm{G})\) 的同态，它是 \(\varGamma(\mathrm{T})\) 到 \(\pi_{1}(\mathrm{T})\) 的典则同构（第 2 节，注 3）与同态 \(\pi_{1}(\iota)\) 的复合，其中 \(\iota\) 是典则单射 \(T\to G\)。

命题 11. 同态 \(f(\mathrm{G},\mathrm{T}) : \Gamma(\mathrm{T}) \to \pi_1(\mathrm{G})\) 是满射。其核是 \(\mathrm{N}(\mathrm{G},\mathrm{T})\)（\(\Gamma(\mathrm{T})\) 的子群），由结点向量族 \((K_{\alpha})_{\alpha \in \mathrm{R}(\mathrm{G},\mathrm{T})}\) 生成。

同态 \(f(\mathrm{G},\mathrm{T})\) 的满射性由 \(\S2\) 第 4 节命题 3 得出。我们用 \(A(\mathrm{G},\mathrm{T})\) 表示断言：「\(f(\mathrm{G},\mathrm{T})\) 的核由诸 \(K_{\alpha}\) 生成」，它尚待证明；我们区分几种情形：

\begin{enumerate}
\def\labelenumi{\alph{enumi})}
\item
  G 是单连通的。设 \(\rho : \mathfrak{g}_{\mathbf{C}} \to \mathfrak{gl}(\mathrm{V})\) 是 \(\mathfrak{g}_{\mathbf{C}}\) 在有限维复向量空间 V 上的一个线性表示。限制到 \(\mathfrak{g}\)，就得到 \(\mathfrak{g}\) 在实向量空间 \(V_{(\mathbf{R})}\) 上的一个表示；由于 G 单连通，存在解析线性表示 \(\pi\)（G 在 \(V_{(\mathbf{R})}\) 上的）使得 \(\rho = L(\pi)\)。由第 3 节命题 7 可知，\(\delta(X(T))\)（X(T) 在 \(\mathfrak{t}_{\mathbf{C}}^{*}\) 中的像）包含 \(\rho\) 在 V 上的所有权。由于这对每个表示 \(\rho\)（\(\mathfrak{g}_{\mathbf{C}}\) 的表示）都成立，由第八章 §7 第 2 节定理 1 可知，\(\delta(X(T))\) 包含 \(\delta(R)\) 的权群，后者按定义是这样的集合：由 \(\lambda \in \mathfrak{t}_{\mathbf{C}}^{*}\) 中满足 \(\lambda(H_{\delta(\alpha)}) \in \mathbf{Z}\)（对所有 \(\alpha \in R\)）者组成，换言之，由满足 \(\lambda(K_{\alpha}) \in 2\pi i\mathbf{Z}\)（对所有 \(\alpha \in R\)）者组成。于是，群 X(T) 包含每个这样的元素 \(\lambda\)：它属于 X(T) \(\otimes\) Q，且 \(\langle\lambda, K_{\alpha} \rangle \in \mathbf{Z}\) 对所有 \(\alpha \in R\) 成立；由对偶性，这蕴含 \(\Gamma(T)\) 由诸 \(K_{\alpha}\) 生成，从而得到断言 \(A(G,T)\)。
\item
  G 是一个单连通群 \(G'\) 与一个环面 S 的直积。此时 T 是 \(T'\)（\(G'\) 的一个极大环面）与 S 的直积，\(\varGamma(\mathrm{T})\) 可等同于 \(\varGamma(\mathrm{T}')\times\varGamma(\mathrm{S})\)，\(\pi_{1}(\mathrm{G})\) 可等同于 \(\pi_{1}(\mathrm{G}')\times\pi_{1}(\mathrm{S})\)，而 \(f(\mathrm{G},\mathrm{T})\) 可等同于分量为 \(f(\mathrm{G}^{\prime},\mathrm{T}^{\prime})\) 与 \(f(\mathrm{S},\mathrm{S})\) 的同态。由于 \(f(\mathrm{S},\mathrm{S})\) 是双射，典则映射 \(\varGamma(\mathrm{T}^{\prime})\to\varGamma(\mathrm{T})\) 把 \(\operatorname{Ker}f(\mathrm{G}^{\prime},\mathrm{T}^{\prime})\) 双射地映到 \(\operatorname{Ker}f(\mathrm{G},\mathrm{T})\) 上。此外，诸 \(K_{\alpha}\) 属于 G 的换位子群 \(G^{\prime}\) 的李代数，从而属于 \(\varGamma(\mathrm{T}^{\prime})\) 的像，故显然 \(A(\mathrm{G}^{\prime},\mathrm{T}^{\prime})\) 蕴含 \(A(\mathrm{G},\mathrm{T})\)，结合 a) 即得断言 \(A(\mathrm{G},\mathrm{T})\)。
\item
  一般情形。存在核为有限的满射态射 \(p : G' \to G\)，其中 \(G'\) 是一个单连通群与一个环面的直积（\(\S1\)，第 4 节，命题 4）。若 \(T'\) 是 T 在 \(G'\) 中的逆像（这是 \(G'\) 的一个极大环面，见 \(\S2\)，第 3 节，命题 1），且 N 是 p 的核，则我们有正合序列 \(0 \to N \to G' \to G \to 0\) 与 \(0 \to N \to T' \to T \to 0\)，从而有一个各行均正合的交换图（第 2 节，注 1 及《一般拓扑》，第十一章，编写中）
\end{enumerate}

\[
\begin{array}{c c c c c c c c c} 0 & \longrightarrow & \Gamma (\mathrm{T} ^ {\prime}) & \longrightarrow & \Gamma (\mathrm{T}) & \longrightarrow & \mathrm{N} & \longrightarrow & 0 \\ & & \Big \downarrow f (\mathrm{G} ^ {\prime}, \mathrm{T} ^ {\prime}) & & \Big \downarrow f (\mathrm{G}, \mathrm{T}) & & \Big \downarrow \mathrm{Id} _ {\mathrm{N}} \\ 0 & \longrightarrow & \pi_ {1} (\mathrm{G} ^ {\prime}) & \longrightarrow & \pi_ {1} (\mathrm{G}) & \longrightarrow & \mathrm{N} & \longrightarrow & 0. \end{array}
\]

由蛇形图（《代数》，第 X 章，第 4 页，命题 2）立即可知，\(A(\mathrm{G}', \mathrm{T}')\) 蕴含 \(A(\mathrm{G}, \mathrm{T})\)，结合 \(b\)) 即得命题。

推论 1. G 是单连通的当且仅当族 \((K_{\alpha})_{\alpha \in \mathrm{R}(\mathrm{G},\mathrm{T})}\) 生成 \(\Gamma (\mathrm{T})\)。

推论 2. 设 H 是 G 的包含 T 的连通闭子群；则有正合序列

\[
0 \longrightarrow \mathrm{N} (\mathrm{H}, \mathrm{T}) \longrightarrow \mathrm{N} (\mathrm{G}, \mathrm{T}) \longrightarrow \pi_ {1} (\mathrm{H}) \longrightarrow \pi_ {1} (\mathrm{G}) \longrightarrow 0.
\]

这由《代数》第 X 章第 4 页命题 2（蛇形图）应用于下面的交换图得出

\[
\begin{array}{c c c c c c c c c} 0 & \longrightarrow & \mathrm{N(H,T)} & \longrightarrow & \varGamma (\mathrm{T}) & \longrightarrow & \pi_ {1} (\mathrm{H}) & \longrightarrow & 0 \\ & & \Big \downarrow & & \Big \downarrow & & \Big \downarrow & & \\ 0 & \longrightarrow & \mathrm{N(G,T)} & \longrightarrow & \varGamma (\mathrm{T}) & \longrightarrow & \pi_ {1} (\mathrm{G}) & \longrightarrow & 0. \end{array}
\]

注. 可以证明（参看 §5 习题 2）\(\pi_2(\mathrm{G / H})\) 为零。于是由前述序列的正合性，得到从 \(\pi_2(\mathrm{G / H})\) 到 \(\mathrm{N(G,T) / N(H,T)}\) 的一个同构。

推论 3. 同态 \(\pi_1(\mathrm{D}(\mathrm{G})) \to \pi_1(\mathrm{G})\)（对应于 \(\mathrm{D}(\mathrm{G})\) 到 \(\mathrm{G}\) 中的包含关系）诱导出从 \(\pi_1(\mathrm{D}(\mathrm{G}))\) 到 \(\pi_1(\mathrm{G})\) 的挠子群的同构。

实际上，\(\mathrm{T} \cap \mathrm{D}(\mathrm{G})\) 是 \(\mathrm{D}(\mathrm{G})\) 的一个极大环面（\(\S 2\)，第 3 节，命题 1 \(c\)）；由正合序列

\[
0 \longrightarrow \Gamma (\mathrm{T} \cap \mathrm{D} (\mathrm{G})) \longrightarrow \Gamma (\mathrm{T}) \longrightarrow \Gamma (\mathrm{T} / (\mathrm{T} \cap \mathrm{D} (\mathrm{G}))) \longrightarrow 0
\]

与命题 11，我们得到正合序列

\[
0 \longrightarrow \pi_ {1} (\mathrm{D} (\mathrm{G})) \longrightarrow \pi_ {1} (\mathrm{G}) \longrightarrow \Gamma (\mathrm{T} / (\mathrm{T} \cap \mathrm{D} (\mathrm{G}))) \longrightarrow 0,
\]

从而得到推论，因为 \(\pi_{1}(\mathrm{D}(\mathrm{G}))\) 有限而 \(\Gamma(\mathrm{T}/(\mathrm{T}\cap\mathrm{D}(\mathrm{G})))\) 自由。

